% method.tex -- Stochastic Attention via Langevin Dynamics on the Hopfield Energy

\subsection{Stochastic attention for energy-based targets}

Let $C:\mathbb R^d\rightarrow\mathbb R$ be a differentiable constraint,
penalty, negative log-likelihood, or other application-specific energy, and let
$\lambda\geq0$. We assume that the combined energy is bounded below strongly
enough to define a normalizable density and that its gradient satisfies the
regularity conditions required by the selected Langevin integrator. The
constrained stochastic-attention target is defined as follows:
\begin{equation}
p_{\beta,\lambda}(\boldsymbol\xi)
\propto \exp\!\left\{-\beta\left[
E_{\beta,\mathbf r}(\boldsymbol\xi)+\lambda C(\boldsymbol\xi)
\right]\right\}.
\label{eq:constrained-target}
\end{equation}
The base stochastic-attention model is the case $\lambda=0$. Applying the
unadjusted Langevin algorithm
to the potential $U_{\beta,\lambda}=\beta(E_{\beta,\mathbf r}+\lambda C)$
with integration step $\eta=\alpha/\beta$ gives the update:
\begin{align}
\boldsymbol\xi_{t+1}
={}&\boldsymbol\xi_t-\alpha\left[
\boldsymbol\xi_t-\mathbf T_{\beta,\mathbf r}(\boldsymbol\xi_t)
+\lambda\nabla C(\boldsymbol\xi_t)\right]
+\sqrt{\frac{2\alpha}{\beta}}\,\boldsymbol\epsilon_t,
\qquad \boldsymbol\epsilon_t\sim\mathcal N(\mathbf0,\mathbf I),
\label{eq:stochastic-attention}\\
={}&(1-\alpha)\boldsymbol\xi_t
+\alpha\mathbf X\operatorname{softmax}\!\left(
\beta\mathbf X^\top\boldsymbol\xi_t+\log\mathbf r\right)
-\alpha\lambda\nabla C(\boldsymbol\xi_t)
+\sqrt{\frac{2\alpha}{\beta}}\,\boldsymbol\epsilon_t.
\label{eq:stochastic-attention-expanded}
\end{align}
This is the stochastic-attention update: an attention-weighted memory drift,
an optional constraint force, and thermally calibrated noise. It requires only
energy gradients and therefore remains applicable when the normalizing
constant and an independent sampler are unavailable. Each base attention step
costs $\mathcal O(dK)$, plus the cost of evaluating $\nabla C$.

\begin{algorithm}[H]
\caption{Stochastic Attention Sampler (ULA)}
\label{alg:stochastic-attention}
\begin{algorithmic}[1]
\REQUIRE Memories $\mathbf X$, inverse temperature $\beta>0$, multiplicities
$\mathbf r\geq\mathbf0$ with $\sum_k r_k>0$, step size $\alpha>0$, iterations
$T$, initial state $\boldsymbol\xi_0$, differentiable energy $C$, weight
$\lambda\geq0$, and a proper target in Eq.~\eqref{eq:constrained-target}
\FOR{$t=0,\ldots,T-1$}
\STATE $\mathbf a_t\leftarrow
\operatorname{softmax}(\beta\mathbf X^\top\boldsymbol\xi_t+\log\mathbf r)$
\STATE $\boldsymbol\epsilon_t\sim\mathcal N(\mathbf0,\mathbf I)$
\STATE $\boldsymbol\xi_{t+1}\leftarrow
(1-\alpha)\boldsymbol\xi_t+\alpha\mathbf X\mathbf a_t
-\alpha\lambda\nabla C(\boldsymbol\xi_t)
+\sqrt{2\alpha/\beta}\,\boldsymbol\epsilon_t$
\ENDFOR
\RETURN $\boldsymbol\xi_0,\ldots,\boldsymbol\xi_T$
\end{algorithmic}
\end{algorithm}

Algorithm~\ref{alg:stochastic-attention} is the general stochastic-attention
sampler. Setting $C\equiv0$ gives the base sampler. Smooth observations,
physical penalties, design objectives, and products of energy experts can be
introduced through $C$. Hard feasible sets require a
constraint-aware variant such as reflected, projected, or proximal Langevin
dynamics; Eq.~\eqref{eq:constrained-target} does not by itself justify applying
an ordinary gradient step to a discontinuous indicator.

\subsection{Regularity and convergence}

Several familiar regimes are recovered directly from
Eq.~\eqref{eq:stochastic-attention}. With $\lambda=0$ and the noise suppressed,
the method is relaxed deterministic attention; setting $\alpha=1$ recovers the
standard Hopfield retrieval map. As $\beta\rightarrow\infty$, a unique maximum
logit selects the maximum-inner-product memory. This is also the nearest memory
when all memories have equal norm. At finite $\beta$, the calibrated noise
produces stochastic attention. Define the local deterministic direction as
$\mathbf d(\boldsymbol\xi)=\mathbf T_{\beta,\mathbf r}(\boldsymbol\xi)
-\boldsymbol\xi-\lambda\nabla C(\boldsymbol\xi)$. The ratio of the drift norm
to the root-mean-square norm of the $d$-dimensional noise increment is:
\begin{equation}
\operatorname{SNR}_{\mathrm{step}}(\boldsymbol\xi)
:=\frac{\alpha\lVert\mathbf d(\boldsymbol\xi)\rVert_2}
{\sqrt{2\alpha d/\beta}}
=\sqrt{\frac{\alpha\beta}{2d}}
\lVert\mathbf d(\boldsymbol\xi)\rVert_2.
\label{eq:snr}
\end{equation}
The commonly reported factor $\sqrt{\alpha\beta/(2d)}$ assumes a unit-scale
drift. It is a scale diagnostic, not a state-independent signal-to-noise ratio
or a universal phase boundary.

The following proposition (proved in
Appendix~\ref{app:proof-regularity}) makes the regularity of the unmodified
Hopfield potential explicit.

\begin{proposition}[Regularity and a sufficient convexity condition]
\label{prop:regularity}
Let $\sigma_{\max}=\lVert\mathbf X\rVert_{\mathrm{op}}$ and
$M=\max_k\lVert\mathbf m_k\rVert_2$. The Hopfield energy has an
$L$-Lipschitz gradient with
$L\leq 1+\beta\sigma_{\max}^2/2$ and satisfies the dissipativity bound:
\begin{equation}
\langle\nabla E(\boldsymbol\xi),\boldsymbol\xi\rangle
\geq \tfrac12\lVert\boldsymbol\xi\rVert_2^2-\tfrac12M^2.
\end{equation}
Moreover, $\beta\sigma_{\max}^2<2$ is a sufficient condition for strict
convexity. It is not claimed to be the exact onset of nonconvexity or a storage
capacity boundary.
\end{proposition}

\begin{corollary}[ULA in the sufficient strongly convex regime]
\label{cor:convergence}
If $\beta\sigma_{\max}^2<2$, then $U=\beta E$ is
$m$-strongly convex with
$m=\beta(1-\beta\sigma_{\max}^2/2)$ and has
$L_U$-Lipschitz gradient with
$L_U=\beta(1+\beta\sigma_{\max}^2/2)$. For a sufficiently small fixed step,
the ULA iterates contract geometrically in $W_2$ toward a step-size-dependent invariant
law, with a discretization error that vanishes as $\alpha\rightarrow0$
\cite{dalalyanTheoreticalGuaranteesApproximate2017,durmusNonasymptoticAnalysisUnadjusted2017}.
\end{corollary}

Outside this sufficient convex regime, the target can be multimodal and
inter-basin mixing can be slow. Warm starts, multiple chains, tempering, and
Metropolis correction address different aspects of finite-time sampling. The
single-chain and temperature experiments in the appendix show how the tested
finite-time trajectories depended on these choices.

The Metropolis-adjusted Langevin algorithm uses the following proposal mean:
\begin{equation}
\boldsymbol\mu(\boldsymbol\xi)
=\boldsymbol\xi-\alpha\left[
\nabla E_{\beta,\mathbf r}(\boldsymbol\xi)
+\lambda\nabla C(\boldsymbol\xi)\right].
\end{equation}
The proposal covariance is $2\alpha\beta^{-1}\mathbf I$. The algorithm accepts
a proposed state $\boldsymbol\xi'$ with probability:
\begin{equation}
\min\!\left\{1,
\frac{p_{\beta,\lambda}(\boldsymbol\xi')
q(\boldsymbol\xi\mid\boldsymbol\xi')}
{p_{\beta,\lambda}(\boldsymbol\xi)
q(\boldsymbol\xi'\mid\boldsymbol\xi)}\right\}.
\label{eq:mala-acceptance}
\end{equation}
The correction preserves the desired target at finite step size. Under the
usual irreducibility and recurrence conditions, the chain converges to that
target. A high acceptance rate does not guarantee rapid mixing and can coexist
with highly correlated moves. Full pseudocode and the proposal-density
correction are given in Appendix~\ref{app:mala}.

\subsection{Exact Gaussian-mixture special case}

For $\lambda=0$, the tied Hopfield target has additional algebraic structure.
Let $p_{\beta,\mathbf r}(\boldsymbol\xi)$ be the normalized Boltzmann density of
Eq.~\eqref{eq:modern-hopfield-energy}.

\begin{theorem}[Exact Gaussian-mixture representation]
\label{thm:exact-mixture}
For $\beta>0$ and at least one positive $r_k$, the normalized density is:
\begin{equation}
p_{\beta,\mathbf r}(\boldsymbol\xi)
=\sum_{k=1}^K w_k\,
\mathcal N(\boldsymbol\xi;\mathbf m_k,\beta^{-1}\mathbf I),
\qquad
w_k=\frac{r_k\exp(\frac\beta2\lVert\mathbf m_k\rVert_2^2)}
{\sum_{j=1}^K r_j\exp(\frac\beta2\lVert\mathbf m_j\rVert_2^2)}.
\label{eq:exact-mixture}
\end{equation}
\end{theorem}

\begin{proof}
Completing the square in the unnormalized density gives the expression:
\begin{align}
\exp[-\beta E_{\beta,\mathbf r}(\boldsymbol\xi)]
&=\sum_{k=1}^K r_k
\exp\!\left(-\frac\beta2\lVert\boldsymbol\xi\rVert_2^2
+\beta\mathbf m_k^\top\boldsymbol\xi\right)\\
&=\sum_{k=1}^K r_k
\exp\!\left(\frac\beta2\lVert\mathbf m_k\rVert_2^2\right)
\exp\!\left[-\frac\beta2
\lVert\boldsymbol\xi-\mathbf m_k\rVert_2^2\right].
\end{align}
Normalizing the common Gaussian factor and the component coefficients yields
Eq.~\eqref{eq:exact-mixture}.
\end{proof}

If the memories have equal norm, including the $\ell_2$-normalized memories in
our experiments, the norm factors cancel; equal multiplicities then give
$w_k=1/K$. The target is an isotropic Parzen estimator with bandwidth
$h=\beta^{-1/2}$. If memory norms differ, uniform component selection is
incorrect and the norm factor in Eq.~\eqref{eq:exact-mixture} is required.

Theorem~\ref{thm:exact-mixture} establishes a precise equivalence for the base
case: continuous-time Langevin dynamics, correctly adjusted MALA, and exact
ancestral sampling share the same target distribution. Finite-step ULA
approximates that target and can be calibrated against exact draws. Exact
ancestral sampling is the efficient method for independent endpoints when
$\lambda=0$. A generic nonlinear $C$ changes the component conditionals and
usually removes this direct sampler, while Algorithms~\ref{alg:stochastic-attention}
and~\ref{alg:mala} continue to use the available constrained score. Constant,
affine, and some quadratic choices of $C$ retain additional analytic structure
and should be simplified directly when possible.

\begin{algorithm}[H]
\caption{Exact sampler for the unmodified tied Hopfield target}
\label{alg:exact-sampler}
\begin{algorithmic}[1]
\REQUIRE Memories $\mathbf X$, inverse temperature $\beta>0$, multiplicities
$\mathbf r\geq\mathbf0$ with $\sum_k r_k>0$, number of endpoints $S$
\STATE $\ell_k\leftarrow
\log r_k+\frac\beta2\lVert\mathbf m_k\rVert_2^2$ for $k=1,\ldots,K$,
with $\log 0=-\infty$
\STATE $\mathbf w\leftarrow\operatorname{softmax}(\boldsymbol\ell)$
\FOR{$s=1,\ldots,S$}
\STATE $z_s\sim\operatorname{Categorical}(\mathbf w)$
\STATE $\boldsymbol\epsilon_s\sim\mathcal N(\mathbf0,\mathbf I)$
\STATE $\boldsymbol\xi_s\leftarrow
\mathbf m_{z_s}+\beta^{-1/2}\boldsymbol\epsilon_s$
\ENDFOR
\RETURN $\{\boldsymbol\xi_s\}_{s=1}^{S}$
\end{algorithmic}
\end{algorithm}

\subsection{Attention as posterior responsibility and score}

Let $Z$ denote the Gaussian component in the base case. Bayes' rule and
Eq.~\eqref{eq:exact-mixture} give its posterior responsibility:
\begin{equation}
\Pr(Z=k\mid\boldsymbol\xi)
=\frac{r_k\exp(\beta\mathbf m_k^\top\boldsymbol\xi)}
{\sum_j r_j\exp(\beta\mathbf m_j^\top\boldsymbol\xi)}.
\label{eq:responsibilities}
\end{equation}
Thus attention weights are the component responsibilities and the attention
output is the posterior mean of the component center. The corresponding score
is:
\begin{equation}
\nabla_{\boldsymbol\xi}\log p_{\beta,\mathbf r}(\boldsymbol\xi)
=\beta\left[\mathbf T_{\beta,\mathbf r}(\boldsymbol\xi)
-\boldsymbol\xi\right].
\label{eq:exact-score}
\end{equation}
For the constrained target, the score becomes
$\beta[\mathbf T_{\beta,\mathbf r}(\boldsymbol\xi)-\boldsymbol\xi
-\lambda\nabla C(\boldsymbol\xi)]$, which is precisely the drift used by
Algorithm~\ref{alg:stochastic-attention}.

\subsection{Responsibility concentration and temperature}

For fixed $\boldsymbol\xi$, write $e_k=\mathbf m_k^\top\boldsymbol\xi$ and
$a_k\propto r_k\exp(\beta e_k)$.

\begin{proposition}[Fixed-query entropy monotonicity]
\label{prop:entropy-derivative}
For equal positive multiplicities, holding $\boldsymbol\xi$ fixed, the
responsibility entropy $H=-\sum_k a_k\log a_k$ satisfies the derivative:
\begin{equation}
\frac{dH}{d\beta}=-\beta\operatorname{Var}_{\mathbf a}(e)\leq0.
\label{eq:entropy-derivative}
\end{equation}
\end{proposition}

\begin{proof}
For equal multiplicities, write $\log a_k=\beta e_k-\log Z$ up to a
constant. Then $H=-\beta\mathbb E_{\mathbf a}[e]+\log Z$. Differentiation,
using $d\log Z/d\beta=\mathbb E_{\mathbf a}[e]$ and
$d\mathbb E_{\mathbf a}[e]/d\beta=\operatorname{Var}_{\mathbf a}(e)$, gives
the result. For unequal fixed multiplicities, an additional covariance term
involving $\log r_k$ appears; throughout the entropy sweeps $r_k=1$.
\end{proof}

For equal multiplicities, entropy decreases monotonically with inverse
temperature for a fixed query. Let
$\mu_3=\mathbb E_{\mathbf a}[(e-\mathbb E_{\mathbf a}[e])^3]$ denote the third
central moment of the logits. The point of maximal entropy loss satisfies the
following inflection condition:
\begin{equation}
\operatorname{Var}_{\mathbf a}(e)\big|_{\beta^*}
=-\beta^*\mu_3\big|_{\beta^*}.
\label{eq:entropy-inflection}
\end{equation}
For random unit memories, inner products have variance of order $1/d$, so the
logit spread becomes order one when $\beta$ is of order $\sqrt d$. This scaling
is an ensemble heuristic, not a universal temperature rule. On a fixed memory
bank, Eq.~\eqref{eq:entropy-inflection} defines a data-dependent descriptive
crossover of the responsibility vector. It does not identify a storage-capacity
boundary or prove a change in sample quality. Fig.~\ref{fig:experiments}
shows the finite-memory concentration diagnostic and a finite-chain
implementation check. Table~\ref{tab:exact-validation} supplies the analytic
equilibrium calibration.

Temperature also has a direct probabilistic interpretation in the base case:
$\beta^{-1/2}$ is the component standard deviation. For unit-norm memories and
an exact draw $\boldsymbol\xi=\mathbf m_k+\boldsymbol\epsilon$, high-dimensional
norm concentration gives the approximation:
\begin{equation}
1-\cos(\boldsymbol\xi,\mathbf m_k)
\approx 1-\frac{1}{\sqrt{1+d/\beta}}.
\label{eq:novelty-approximation}
\end{equation}
This approximation neglects the noise projection along $\mathbf m_k$ and
assumes the generating component remains the nearest memory.

\subsection{Masking and multiplicity}

For a selected memory subset $\mathcal S$, define
$b_k=0$ for $k\in\mathcal S$ and $b_k=-\infty$ otherwise. The base masked
stochastic-attention update is:
\begin{equation}
\boldsymbol\xi_{t+1}
=(1-\alpha)\boldsymbol\xi_t
+\alpha\mathbf X\operatorname{softmax}\!\left(
\beta\mathbf X^\top\boldsymbol\xi_t+\mathbf b\right)
+\sqrt{\frac{2\alpha}{\beta}}\,\boldsymbol\epsilon_t.
\label{eq:masked-sa}
\end{equation}
It is identical to running stochastic attention on the reduced memory matrix
$\mathbf X_{\mathcal S}$. Equivalently, setting $r_k=0$ removes component $k$;
positive multiplicities bias logits by $\log r_k$, and finite biases satisfy
$r_k=e^{b_k}$. In the unmodified tied model, these operations also reweight or
restrict the exact mixture. In a constrained target they remain part of the
stochastic-attention energy and can be combined with $C$. This provides class-
or subset-conditioned sampling without retraining while keeping the distinction
between a memory mask and a general state-dependent constraint explicit.
